\section{The Greatest Common Divisor (GCD)}

\begin{definition}[Greatest Common Divisor]
Let $a, b \in \mathbb{Z}$, such that $a \neq 0$ or $b \neq 0$. The set of their common divisors has a largest element $d \in \mathbb{N}$, called the greatest common divisor of $a$ and $b$. We write $d = \gcd(a, b)$.
\end{definition}

\bigskip

\begin{remark}[Example of GCD]
To find $\gcd(18, 12)$, we can list the divisors of each number:
\begin{itemize}
    \item Divisors of 12: $\{1, 2, 3, 4, 6, 12\}$
    \item Divisors of 18: $\{1, 2, 3, 6, 9, 18\}$
\end{itemize}
The set of common divisors is $\{1, 2, 3, 6\}$. The largest element in this set is $6$, so $\gcd(18, 12) = 6$.
\par
While listing all divisors is feasible for small numbers, it is impractical for large integers. A far more powerful and efficient method is the Euclidean Algorithm.
\end{remark}
